\chapter{Combinatorics and Probability}

\begin{theorem}[Addition and Multiplication Principles]
If mutually exclusive classes of objects have respective cardinalities $a_1, \dots, a_k$, their union has cardinality $a_1+\cdots+a_k$. If an object is determined by $k$ independent choices with $a_i$ possibilities at step $i$, the number of objects is $a_1\cdots a_k$.
\end{theorem}

\begin{definition}[Factorials, Permutations, and Combinations]
For $n\in\Nzero$, let $n!=\prod_{k=1}^n k$, with $0!=1$. If $0\le r\le n$, the number of ordered selections of $r$ distinct elements from an $n$-element set is
\[
 {}_nP_r=\frac{n!}{(n-r)!},
\]
and the number of $r$-element subsets is
\[
 \binom nr=\frac{n!}{r!(n-r)!}.
\]
\end{definition}

\begin{formula}[Multiset Combinations and Stars and Bars]
For $n\ge1$ and $r\ge0$, the number of multisets of cardinality $r$ chosen from $n$ types is
\[
 \multichoose nr=\binom{n+r-1}r.
\]
Equivalently, this is the number of nonnegative integer solutions of $x_1+\cdots+x_n=r$. The number of positive integer solutions for $r\ge n$ is $\binom{r-1}{n-1}$.
\end{formula}

\begin{theorem}[Binomial and Multinomial Theorems]
For $n\in\Nzero$,
\[
 (x+y)^n=\sum_{k=0}^n\binom nk x^ky^{n-k}.
\]
If $n_1, \dots, n_m\in\Nzero$ sum to $n$, the corresponding multinomial coefficient is
\[
 \binom{n}{n_1, \dots, n_m}=\frac{n!}{n_1!\cdots n_m!},
\]
and $(x_1+\cdots+x_m)^n$ expands as the sum of these coefficients times $x_1^{n_1}\cdots x_m^{n_m}$.
\end{theorem}

\begin{formula}[Basic Binomial Identities]
For $0\le k\le n$,
\[
 \binom nk=\binom n{n-k}, \qquad
 \binom nk+\binom n{k-1}=\binom{n+1}k,
\]
where coefficients with inadmissible lower indices are zero. Furthermore,
\[
 \sum_{k=0}^n\binom nk=2^n, \quad
 \sum_{k=0}^n(-1)^k\binom nk=0\quad(n\ge1),
\]
\[
 \sum_{k=0}^n k\binom nk=n2^{n-1}\quad(n\ge1).
\]
\end{formula}

\begin{theorem}[Vandermonde's Identity]
For nonnegative integers $r, s, n$,
\[
 \sum_{k=0}^n\binom rk\binom s{n-k}=\binom{r+s}n.
\]
In particular,
\[
 \sum_{k=0}^n\binom nk^2=\binom{2n}n.
\]
\end{theorem}

\begin{theorem}[Binomial Inversion]
For sequences $(a_n)$ and $(b_n)$ over a field of characteristic zero,
\[
 b_n=\sum_{k=0}^n\binom nk a_k
 \quad\Longleftrightarrow\quad
 a_n=\sum_{k=0}^n(-1)^{n-k}\binom nk b_k.
\]
\end{theorem}

\begin{theorem}[Cauchy--Davenport Theorem]
Let $A,B$ be nonempty subsets of the cyclic group $\Z/p\Z$ for a prime $p$. Their sumset $A+B=\{a+b\mid a\in A,\ b\in B\}$ satisfies
\[
 |A+B|\ge\min(p,\,|A|+|B|-1).
\]
\end{theorem}

\begin{theorem}[Pigeonhole Principle]
If $N$ objects are placed in $k\ge1$ boxes, some box contains at least $\lceil N/k\rceil$ objects. In particular, placing $k+1$ objects in $k$ boxes forces at least two into the same box.
\end{theorem}

\begin{theorem}[Inclusion--Exclusion Principle]
For finite sets $A_1, \dots, A_n$,
\[
 \left|\bigcup_{i=1}^n A_i\right|
 =\sum_{\varnothing\ne S\subseteq[n]}(-1)^{|S|+1}
 \left|\bigcap_{i\in S}A_i\right|.
\]
\end{theorem}

\begin{formula}[Counting Surjections]
For $n\ge0$ and $k\ge1$, the number of surjections from an $n$-element set to a $k$-element set is
\[
 \sum_{j=0}^k(-1)^j\binom kj(k-j)^n=k!\, S(n, k),
\]
where $S(n, k)$ is the Stirling number of the second kind and $0^0=1$ in the finite counting sum.
\end{formula}

\begin{definition}[Derangement]
A derangement is a permutation with no fixed point. Write $D_n$ or $!n$ for the number of derangements of $n$ elements; here $D_n$ is not the dihedral group, the notation being determined by context.
\end{definition}

\begin{formula}[Derangement Formula]
For $n\in\Nzero$,
\[
 D_n=n!\sum_{k=0}^n\frac{(-1)^k}{k!}.
\]
For $n\ge2$, $D_n=(n-1)(D_{n-1}+D_{n-2})$, with $D_0=1$ and $D_1=0$. For $n\ge1$, $D_n$ is the nearest integer to $n!/e$.
\end{formula}

\begin{definition}[Rook Polynomial]
For a finite set $B$ of forbidden cells on an $n\times n$ board, let $r_k(B)$ count placements of $k$ pairwise nonattacking rooks entirely in $B$, with $r_0(B)=1$. The rook polynomial is
\[
 R_B(x)=\sum_{k=0}^n r_k(B)x^k.
\]
\end{definition}

\begin{formula}[Forbidden Positions and Rook Numbers]
The number of permutations $\sigma\in S_n$ whose selected cells $(i,\sigma(i))$ avoid a set $B$ of forbidden cells is
\[
 \sum_{k=0}^n(-1)^k r_k(B)(n-k)!.
\]
For the diagonal forbidden board, $r_k(B)=\binom nk$, giving the derangement formula.
\end{formula}

\begin{definition}[Stirling Numbers and Bell Numbers]
The unsigned Stirling number $c(n, k)$ counts permutations of $[n]$ with $k$ cycles. The Stirling number $S(n, k)$ counts partitions of $[n]$ into $k$ nonempty blocks. The Bell number
\[
 B_n=\sum_{k=0}^n S(n, k)
\]
counts all set partitions of $[n]$.
\end{definition}

\begin{formula}[Stirling Number Recurrences]
For $n, k\ge1$,
\[
 c(n, k)=c(n-1, k-1)+(n-1)c(n-1, k),
\]
\[
 S(n, k)=S(n-1, k-1)+kS(n-1, k).
\]
The boundary conditions are $c(0, 0)=S(0, 0)=1$ and zero for inadmissible indices.
\end{formula}

\begin{definition}[Latin Squares]
A Latin square of order $n$ is an $n\times n$ array of $n$ symbols, each occurring exactly once in each row and each column.
\end{definition}

\begin{definition}[Block Designs and Steiner Triple Systems]
A $2$-$(v,k,\lambda)$ design is a $v$-element set together with a family of $k$-element blocks in which each pair of distinct points is contained in exactly $\lambda$ blocks. A Steiner triple system is a $2$-$(v,3,1)$ design. For $v\ge3$, such a system exists if and only if $v\equiv1$ or $3\pmod6$.
\end{definition}

\begin{definition}[Catalan Numbers]
For $n\in\Nzero$, the $n$th Catalan number is
\[
 C_n=\frac{1}{n+1}\binom{2n}n
 =\binom{2n}n-\binom{2n}{n+1}.
\]
It counts, among other objects, correctly matched parenthesis strings of length $2n$ and lattice paths from $(0, 0)$ to $(n, n)$ that never rise above the diagonal.
\end{definition}

\begin{formula}[Catalan Recurrence and Generating Function]
The Catalan numbers satisfy
\[
 C_0=1, \qquad C_{n+1}=\sum_{k=0}^n C_kC_{n-k}.
\]
Their ordinary generating function is
\[
 C(x)=\sum_{n\ge0}C_nx^n
 =\frac{1-\sqrt{1-4x}}{2x},
\]
with the constant term understood by formal-power-series expansion.
\end{formula}

\begin{definition}[Integer Partitions and Compositions]
A partition of $n\in\Nzero$ is a weakly decreasing sequence of positive integers summing to $n$, with the empty partition for $n=0$. A composition of $n\in\N$ is an ordered sequence of positive integers summing to $n$. The number of compositions of $n$ is $2^{n-1}$.
\end{definition}

\begin{theorem}[Euler's Partition Generating Function]
If $p(n)$ is the number of integer partitions of $n$ with $p(0)=1$, then, as a formal power series,
\[
 \sum_{n\ge0}p(n)x^n=\prod_{k\ge1}\frac1{1-x^k}.
\]
The number of partitions into distinct parts equals the number of partitions into odd parts.
\end{theorem}

\begin{theorem}[Reflection Principle for Lattice Paths]
For nonnegative integers $a, b$ with $a\ge b$, the number of lattice paths from $(0, 0)$ to $(a, b)$ using right and up steps and never going above $y=x$ is
\[
 \binom{a+b}b-\binom{a+b}{b-1}
 =\frac{a-b+1}{a+1}\binom{a+b}b.
\]
For $a=b=n$, this equals $C_n$.
\end{theorem}

\begin{formula}[Ballot Theorem]
Suppose candidate $A$ receives $a$ votes and candidate $B$ receives $b$ votes, with $a>b\ge0$. Among all $\binom{a+b}{a}$ orders of the votes, the fraction for which $A$ is strictly ahead of $B$ after every vote is
\[
 \frac{a-b}{a+b}.
\]
Consequently, the number of such orders is $\frac{a-b}{a+b}\binom{a+b}{a}$.
\end{formula}

\begin{definition}[Recurrence Relation and Generating Function]
A recurrence relation specifies sequence terms in terms of earlier terms. The ordinary generating function of $(a_n)_{n\ge0}$ is
\[
 A(x)=\sum_{n\ge0}a_nx^n,
\]
viewed either formally or as an analytic series where it converges. The exponential generating function is $\sum_{n\ge0}a_nx^n/n!$.
\end{definition}

\begin{formula}[Geometric and Binomial Series]
As formal power series, for $m\in\N$,
\[
 \frac1{1-x}=\sum_{n\ge0}x^n, \qquad
 \frac1{(1-x)^m}=\sum_{n\ge0}\binom{n+m-1}{m-1}x^n.
\]
The same identities hold analytically when $|x|<1$.
\end{formula}

\begin{theorem}[Generating Functions for Linear Recurrences]
A sequence satisfying a linear recurrence with constant coefficients for all sufficiently large indices has a rational ordinary generating function. Conversely, the coefficients of any rational formal power series satisfy a constant-coefficient linear recurrence for all sufficiently large indices.
\end{theorem}

\begin{formula}[Coefficient Extraction and Convolution]
For formal power series $A(x)=\sum_{n\ge0}a_nx^n$ and $B(x)=\sum_{n\ge0}b_nx^n$,
\[
 [x^n](A(x)B(x))=\sum_{k=0}^n a_kb_{n-k}.
\]
More generally, $[x^n]x^kA(x)=a_{n-k}$ for $n\ge k$.
\end{formula}

\begin{formula}[Exponential Generating Functions]
For a sequence $(a_n)_{n\ge0}$, define its exponential generating function by
\[
 A(x)=\sum_{n\ge0}a_n\frac{x^n}{n!}.
\]
If $B(x)=\sum_{n\ge0}b_nx^n/n!$, then
\[
 n![x^n](A(x)B(x))
 =\sum_{k=0}^n\binom nk a_kb_{n-k}.
\]
This is the labeled-product convolution, in contrast with ordinary generating functions.
\end{formula}
\begin{theorem}[Exponential Formula for Labeled Structures]
Let $\mathcal{C}$ be a class of labeled connected structures, with $c_n$ structures on an $n$-element label set. Suppose $\mathcal{A}$ consists of all finite unordered sets of $\mathcal{C}$-structures on pairwise disjoint nonempty label sets, including the empty set, and this decomposition is unique. If $a_n$ counts objects of $\mathcal{A}$ on $[n]$, then, as formal power series,
\[
 \sum_{n\ge0}a_n\frac{x^n}{n!}=
 \exp\left(\sum_{n\ge1}c_n\frac{x^n}{n!}\right).
\]
\end{theorem}

\begin{theorem}[Lagrange Inversion, Basic Form]
Let $\phi(t)$ be a formal power series over a field of characteristic zero with $\phi(0)\ne0$. If $T(x)$ is the unique formal series with $T(0)=0$ satisfying
\[
 T=x\phi(T),
\]
then for $n\ge1$,
\[
 [x^n]T(x)=\frac1n[t^{n-1}]\phi(t)^n.
\]
\end{theorem}

\begin{theorem}[Burnside's Lemma]
If a finite group $G$ acts on a finite set $X$, the number of orbits is
\[
 |G\backslash X|=\frac1{|G|}\sum_{g\in G}|\Fix(g)|,
\]
where $\Fix(g)=\{x\in X\mid g\cdot x=x\}$.
\end{theorem}

\begin{definition}[Cycle Index of a Permutation Group]
Let a finite group $G$ act on $n$ positions. Write $c_j(g)$ for the number of cycles of length $j$ in the permutation induced by $g\in G$. The cycle index polynomial is
\[
 Z_G(z_1,\dots,z_n)
 =\frac1{|G|}\sum_{g\in G}\prod_{j=1}^n z_j^{c_j(g)}.
\]
In particular, $\sum_{j=1}^n j c_j(g)=n$.
\end{definition}

\begin{theorem}[P\'olya Enumeration Theorem]
Let a finite group $G$ permute an $n$-element set, and suppose each position receives one of $q\in\N$ colors. If $c(g)$ denotes the number of cycles of $g$ in its permutation of the positions, the number of colorings up to $G$-symmetry is
\[
 \frac1{|G|}\sum_{g\in G}q^{c(g)}.
\]
\end{theorem}

\begin{proposition}[Weighted P\'olya Enumeration]
Let $G$ act on $n$ positions, with cycle-index polynomial $Z_G$. If the available colors have respective weights $w_1,\dots,w_q$, the generating polynomial of colorings up to $G$-symmetry, weighted by color multiplicities, is
\[
 Z_G\!\left(\sum_{i=1}^q w_i,\ \sum_{i=1}^q w_i^2,\ \dots,\ \sum_{i=1}^q w_i^n\right).
\]
The coefficient of $\prod_iw_i^{m_i}$ counts the orbits in which color $i$ occurs exactly $m_i$ times. Substituting $w_i=1$ yields $Z_G(q,\dots,q)$.
\end{proposition}

\begin{definition}[Graphs and Degrees]
A finite simple undirected graph is $G=(V,E)$, where $E\subseteq\binom V2$. Its vertex degree $\deg(v)$ is the number of incident edges. A subgraph is determined by subsets of the vertices and edges; an induced subgraph on $S\subseteq V$ contains all edges of $G$ with both endpoints in $S$.
\end{definition}

\begin{definition}[Walks, Trails, Connectivity, and Components]
A walk is a sequence of adjacent vertices, allowing repetition. A trail is a walk with no repeated edges, and a simple path has no repeated vertices. A graph is connected if every pair of vertices can be joined by a path. A connected component is a maximal connected subgraph. A bridge is an edge whose deletion increases the number of connected components; a cut vertex is a vertex whose deletion increases that number.
\end{definition}

\begin{theorem}[Handshaking Lemma]
For any finite undirected graph,
\[
 \sum_{v\in V(G)}\deg(v)=2|E(G)|.
\]
Consequently, the number of vertices of odd degree is even.
\end{theorem}

\begin{theorem}[Erd\H{o}s--Gallai Theorem]
Nonnegative integers $d_1\ge\cdots\ge d_n$ are the degrees of a simple graph on $n$ labeled vertices if and only if $\sum_i d_i$ is even and, for every $1\le k\le n$,
\[
 \sum_{i=1}^k d_i\le k(k-1)+\sum_{i=k+1}^n\min(d_i,k).
\]
\end{theorem}


\begin{definition}[Trees and Spanning Trees]
A tree is a connected graph with no cycles. A spanning tree of a connected graph is a subgraph that is a tree and contains every vertex.
\end{definition}

\begin{theorem}[Equivalent Characterizations of a Tree]
For a finite simple graph $G$ with $n\ge1$ vertices, any two of the following properties imply the third: $G$ is connected; $G$ is acyclic; $G$ has exactly $n-1$ edges. In particular, a tree on $n$ vertices has $n-1$ edges and there is a unique simple path between any two of its vertices.
\end{theorem}

\begin{theorem}[Cayley's Formula]
For $n\ge2$, the number of labeled trees on vertex set $[n]$ is
\[
 n^{n-2}.
\]
For $n=1$, there is exactly one such tree.
\end{theorem}

\begin{theorem}[Pr\"ufer Codes]
Labeled trees on vertex set $[n]$, for $n\ge2$, are in bijection with sequences in $[n]^{n-2}$ (Pr\"ufer codes). The degree of vertex $i$ is one plus its number of appearances in the code.
\end{theorem}

\begin{definition}[Bipartite Graph]
A graph $G=(V, E)$ is bipartite if $V$ can be partitioned into disjoint sets $U, W$ so that every edge joins a vertex of $U$ to a vertex of $W$.
\end{definition}

\begin{theorem}[Bipartite Graph Criterion]
A finite graph is bipartite if and only if it has no cycle of odd length.
\end{theorem}


\begin{definition}[Vertex Coloring and Chromatic Number]
A proper vertex coloring of a graph $G$ assigns colors to the vertices so that adjacent vertices receive distinct colors. The \emph{chromatic number} $\chi(G)$ is the smallest number of colors in a proper coloring.
\end{definition}

\begin{theorem}[Greedy Coloring Bound]
Every finite simple graph $G$ satisfies
\[
 \chi(G)\le\Delta(G)+1,
\]
where $\Delta(G)$ is its maximum vertex degree. Every bipartite graph with at least one edge has chromatic number $2$.
\end{theorem}

\begin{definition}[Chromatic Polynomial]
For a finite graph $G$, the chromatic polynomial $\chi_G(q)$ is the number of proper vertex colorings using $q$ labeled colors, as a polynomial in $q$. It has degree $|V(G)|$ and leading coefficient $1$. If $G$ contains a loop, $\chi_G(q)=0$.
\end{definition}

\begin{theorem}[Deletion--Contraction Formula]
For any nonloop edge $e$ of a finite graph $G$,
\[
 \chi_G(q)=\chi_{G-e}(q)-\chi_{G/e}(q),
\]
where $G-e$ deletes $e$ and $G/e$ contracts $e$. Multiple edges arising under contraction do not change the proper coloring count. For the edgeless graph on $n$ vertices, $\chi_G(q)=q^n$.
\end{theorem}

\begin{theorem}[Dirac's and Ore's Hamiltonian Criteria]
Let $G$ be a finite simple graph with $n\ge3$ vertices. If every vertex has degree at least $n/2$, then $G$ has a Hamiltonian cycle (Dirac). More generally, if $\deg(u)+\deg(v)\ge n$ for every pair of distinct nonadjacent vertices $u, v$, then $G$ has a Hamiltonian cycle (Ore).
\end{theorem}

\begin{theorem}[Eulerian Trail and Circuit Criteria]
A finite connected graph has an Eulerian circuit exactly when every vertex has even degree. It has an Eulerian trail with two distinct endpoints exactly when precisely two vertices have odd degree.
\end{theorem}

\begin{theorem}[Euler's Formula for Planar Graphs]
If a finite connected graph has a plane embedding with $V$ vertices, $E$ edges, and $F$ faces, including the unbounded face, then
\[
 V-E+F=2.
\]
For a simple planar graph with $V\ge3$,
\[
 E\le3V-6.
\]
If it is also bipartite, then $E\le2V-4$.
\end{theorem}

\begin{theorem}[Edge Bounds for Simple Planar Graphs]
A simple planar graph with $v\ge3$ vertices and $e$ edges satisfies $e\le 3v-6$. If it is also bipartite, then $e\le 2v-4$. In particular, every simple planar graph has a vertex of degree at most $5$.
\end{theorem}

\begin{theorem}[Kuratowski's Theorem]
A finite graph is planar if and only if it contains no subdivision of $K_5$ or $K_{3, 3}$ as a subgraph.
\end{theorem}

\begin{definition}[Matching and Vertex Cover]
A matching is a set of edges no two of which share a vertex. A matching is perfect if every vertex is incident with a matching edge. A vertex cover is a set of vertices incident with every edge.
\end{definition}

\begin{theorem}[Hall's Marriage Theorem]
A finite bipartite graph with vertex classes $U, W$ has a matching saturating every vertex of $U$ if and only if
\[
 |N(S)|\ge |S|\qquad\text{for every }S\subseteq U,
\]
where $N(S)$ is the set of neighbors of vertices in $S$.
\end{theorem}

\begin{theorem}[K\H{o}nig's Theorem]
In every finite bipartite graph, the maximum size of a matching equals the minimum size of a vertex cover.
\end{theorem}

\begin{theorem}[Tutte's Perfect Matching Theorem]
A finite undirected graph $G$ has a perfect matching if and only if, for every $S\subseteq V(G)$, the graph $G-S$ has at most $|S|$ connected components of odd order.
\end{theorem}

\begin{theorem}[Stable Marriage Theorem]
For two finite equally sized sets of applicants, each ranking all members of the opposite set strictly, there exists a perfect matching with no blocking pair. A blocking pair consists of two unmatched partners who both prefer one another to their assigned partners.
\end{theorem}

\begin{definition}[Flow Networks and Cuts]
Let $G=(V,E)$ be a finite directed graph with source $s$, sink $t$, and capacities $c(e)\ge0$. A feasible flow assigns $0\le f(e)\le c(e)$ to each edge, with
\[
 \sum_{(u,v)\in E}f(u,v)=\sum_{(v,w)\in E}f(v,w)
 \qquad(v\notin\{s,t\}).
\]
Its value is the net outflow at $s$. An $s$--$t$ cut is a partition $V=S\sqcup T$ with $s\in S$ and $t\in T$; its capacity is $\sum_{u\in S,\,v\in T}c(u,v)$ over edges directed from $S$ to $T$.
\end{definition}

\begin{theorem}[Max-Flow Min-Cut Theorem]
In a finite directed capacitated network with nonnegative edge capacities and distinguished source $s$ and sink $t$, the maximum possible value of an $s$--$t$ flow equals the minimum capacity of an $s$--$t$ cut.
\end{theorem}

\begin{theorem}[Menger's Theorem, Vertex Form]
For distinct nonadjacent vertices $u, v$ in a finite graph, the maximum number of internally vertex-disjoint $u$--$v$ paths equals the minimum number of vertices whose removal separates $u$ and $v$.
\end{theorem}

\begin{theorem}[Matrix--Tree Theorem]
Let $G$ be a finite connected graph with adjacency matrix $\mathbf{A}$ and degree matrix $\mathbf{D}$. Set $\mathbf{L}=\mathbf{D}-\mathbf{A}$. The number of spanning trees of $G$ equals any principal cofactor of $\mathbf{L}$.
\end{theorem}

\begin{definition}[Partially Ordered Sets, Chains, and Antichains]
A partially ordered set (poset) is a set $X$ with a reflexive, antisymmetric, transitive relation $\preceq$. A chain is a subset whose elements are pairwise comparable; an antichain is a subset whose distinct elements are pairwise incomparable.
\end{definition}

\begin{theorem}[Dilworth's Theorem]
In a finite poset, the maximum cardinality of an antichain equals the minimum number of chains in a partition of the poset.
\end{theorem}

\begin{theorem}[Sperner's Theorem]
For a family $\mathcal{A}$ of subsets of $[n]$, no one of which properly contains another,
\[
 |\mathcal{A}|\le\binom n{\lfloor n/2\rfloor}.
\]
\end{theorem}

\begin{theorem}[Lubell--Yamamoto--Meshalkin Inequality]
If $\mathcal{A}$ is an antichain in the Boolean lattice $2^{[n]}$, then
\[
 \sum_{A\in\mathcal{A}}\frac1{\binom n{|A|}}\le1.
\]
\end{theorem}

\begin{definition}[Finite Impartial Games and Nim-Sum]
A finite impartial normal-play game has finitely many positions, no infinite sequence of legal moves, and the same available moves for both players; the player with no legal move loses. The nim-sum $a\oplus b$ is the bitwise exclusive-or of the binary expansions of nonnegative integers $a,b$.
\end{definition}
\begin{theorem}[Sprague--Grundy Theorem]
For a finite impartial normal-play game, define the Grundy value recursively by
\[
 g(P)=\operatorname{mex}\{g(Q)\mid P\to Q\},
\]
where $\operatorname{mex}$ is the least nonnegative integer absent from a set and terminal positions have value $0$. A position is losing for the player to move if and only if $g(P)=0$. The Grundy value of a disjunctive sum of games is the nim-sum of their Grundy values.
\end{theorem}
\begin{theorem}[Bouton's Theorem for Nim]
For a position in ordinary Nim with heaps of sizes $a_1,\dots,a_k\in\Nzero$, in which a move removes a positive number of objects from exactly one heap and the player making the last move wins, the position is losing if and only if
\[
 a_1\oplus\cdots\oplus a_k=0.
\]
\end{theorem}

\begin{theorem}[Ramsey's Theorem]
For positive integers $r, k_1, \dots, k_m$, there exists $N$ such that every $m$-coloring of the $r$-element subsets of $[N]$ contains a set of size $k_i$ all of whose $r$-element subsets have color $i$, for at least one index $i$. In the graph case, $R(s, t)$ is the least $N$ such that every red--blue edge coloring of $K_N$ has a red $K_s$ or a blue $K_t$.
\end{theorem}

\begin{theorem}[Erd\H{o}s--Szekeres Monotone Subsequence Theorem]
Every sequence of $(r-1)(s-1)+1$ distinct real numbers contains an increasing subsequence of length $r$ or a decreasing subsequence of length $s$.
\end{theorem}

\begin{theorem}[Tur\'an's Theorem]
Among simple graphs on $n$ vertices containing no complete graph $K_{r+1}$, the unique graph with the maximum number of edges (up to isomorphism) is the complete $r$-partite graph whose part sizes differ by at most one.
\end{theorem}

\begin{definition}[Probability Space and Conditional Probability]
A probability space is a triple $(\Omega, \mathcal{F}, \Prob)$, where $\mathcal{F}$ is a $\sigma$-algebra of subsets of $\Omega$ and $\Prob$ is a countably additive probability measure with $\Prob(\Omega)=1$. For $\Prob(B)>0$,
\[
 \Prob(A\mid B)=\frac{\Prob(A\cap B)}{\Prob(B)}.
\]
For a finite uniform space $\Omega$, $\Prob(A)=|A|/|\Omega|$.
\end{definition}

\begin{theorem}[Addition Rule and Union Bound]
For events $A, B$,
\[
 \Prob(A\cup B)=\Prob(A)+\Prob(B)-\Prob(A\cap B).
\]
For any finite or countable family of events,
\[
 \Prob\!\left(\bigcup_i A_i\right)\le\sum_i\Prob(A_i).
\]
\end{theorem}

\begin{theorem}[Law of Total Probability and Bayes' Theorem]
Let $(B_i)_{i=1}^m$ be a partition of $\Omega$ into events of positive probability. For any event $A$,
\[
 \Prob(A)=\sum_{i=1}^m\Prob(A\mid B_i)\Prob(B_i).
\]
If $\Prob(A)>0$, then
\[
 \Prob(B_j\mid A)=
 \frac{\Prob(A\mid B_j)\Prob(B_j)}
 {\sum_{i=1}^m\Prob(A\mid B_i)\Prob(B_i)}.
\]
\end{theorem}

\begin{definition}[Independence]
Events $A, B$ are independent when $\Prob(A\cap B)=\Prob(A)\Prob(B)$. A family of events is mutually independent if this product identity holds for every finite nonempty subfamily. Random variables $X, Y$ are independent if the events $\{X\in U\}$ and $\{Y\in V\}$ are independent for all Borel sets $U, V$.
\end{definition}

\begin{definition}[Random Variable and Expectation]
A real-valued random variable $X$ is a measurable map from $\Omega$ to $\R$. Its expectation, when defined, is denoted $\Exp[X]$. For a discrete random variable with countable range and $\sum_x |x|\Prob(X=x)<\infty$,
\[
 \Exp[X]=\sum_x x\Prob(X=x).
\]
\end{definition}


\begin{definition}[Distribution of a Random Variable]
The distribution of $X$ is specified by $\Prob(X\in B)$ for Borel sets $B$. For a discrete random variable, its probability mass function is $p_X(x)=\Prob(X=x)$; its cumulative distribution function is $F_X(x)=\Prob(X\le x)$.
\end{definition}

\begin{definition}[Geometric Probability]
For a bounded measurable planar region $D$ of positive finite area, a point $X$ chosen uniformly from $D$ satisfies
\[
 \Prob(X\in A)=\frac{\operatorname{Area}(A\cap D)}{\operatorname{Area}(D)}
\]
for every measurable $A$. In three dimensions, volume replaces area. The probability model and sampling measure must be specified before applying such ratios.
\end{definition}

\begin{formula}[Bernoulli, Binomial, and Geometric Distributions]
If $X\sim\Bern(p)$, then $\Prob(X=1)=p$ and $\Exp[X]=p$, $\Var(X)=p(1-p)$. If $X\sim\Bin(n, p)$, then for $0\le k\le n$,
\[
 \Prob(X=k)=\binom nk p^k(1-p)^{n-k},
 \quad \Exp[X]=np, \quad\Var(X)=np(1-p).
\]
If $Y\sim\Geom(p)$ counts the number of independent Bernoulli trials up to and including the first success, with $0<p\le1$, then
\[
 \Prob(Y=k)=(1-p)^{k-1}p\quad(k\in\N),
 \qquad \Exp[Y]=\frac1p.
\]
\end{formula}

\begin{formula}[Hypergeometric Distribution]
If a sample of size $n$ is drawn uniformly without replacement from $N$ objects, of which $K$ are designated successes, the number $X$ of successes satisfies
\[
 \Prob(X=k)=\frac{\binom Kk\binom{N-K}{n-k}}{\binom Nn}
\]
for integers $k$ with $\max(0, n-(N-K))\le k\le\min(n, K)$, and $\Exp[X]=nK/N$.
\end{formula}

\begin{formula}[Poisson Distribution and Poisson Limit]
For $\lambda>0$, a Poisson random variable $X\sim\operatorname{Poi}(\lambda)$ has
\[
 \Prob(X=k)=e^{-\lambda}\frac{\lambda^k}{k!},
 \quad k\in\Nzero,\qquad
 \Exp[X]=\Var(X)=\lambda.
\]
If $X_n\sim\Bin(n,p_n)$ with $np_n\to\lambda$, then for every fixed $k\in\Nzero$,
\[
 \Prob(X_n=k)\longrightarrow e^{-\lambda}\frac{\lambda^k}{k!}.
\]
\end{formula}
\begin{formula}[Uniform and Exponential Distributions]
If $X$ is uniform on $[a,b]$ with $a<b$, its density is $f_X(x)=1/(b-a)$ for $a\le x\le b$ (and $0$ otherwise), and $\Exp[X]=(a+b)/2$. If $Y$ is exponential with rate $\lambda>0$, its density is
\[
 f_Y(y)=\lambda e^{-\lambda y}\quad(y\ge0),
 \qquad \Exp[Y]=\frac1\lambda,\quad \Var(Y)=\frac1{\lambda^2}.
\]
\end{formula}

\begin{definition}[Finite Markov Chains]
A finite-state, time-homogeneous Markov chain has a row-stochastic transition matrix $\mathbf{P}=(p_{ij})$ satisfying
\[
 \Prob(X_{n+1}=j\mid X_n=i,\dots,X_0)=p_{ij}.
\]
A probability row vector $\boldsymbol\pi$ is stationary if $\boldsymbol\pi\mathbf{P}=\boldsymbol\pi$.
\end{definition}

\begin{theorem}[Stationary Distributions of Finite Markov Chains]
Every finite Markov chain has a stationary distribution. If its transition graph is irreducible, the stationary distribution is unique and strictly positive. If it is also aperiodic, then $\Prob(X_n=j\mid X_0=i)\to\pi_j$ for every pair of states $i,j$.
\end{theorem}

\begin{definition}[Normal Distribution]
A real random variable $X$ has normal distribution $\Normal(\mu,\sigma^2)$, where $\sigma>0$, if its density is
\[
 f(x)=\frac1{\sigma\sqrt{2\pi}}\exp\!\left(-\frac{(x-\mu)^2}{2\sigma^2}\right),\qquad x\in\R.
\]
Its expectation is $\mu$, its variance is $\sigma^2$, and $(X-\mu)/\sigma\sim\Normal(0,1)$.
\end{definition}

\begin{definition}[Conditional Expectation]
For integrable $X$ and a random variable $Y$, the conditional expectation $\Exp[X\mid Y]$ denotes an integrable function of $Y$ whose expectation on each event determined by $Y$ agrees with that of $X$. In the discrete case, when $\Prob(Y=y)>0$,
\[
 \Exp[X\mid Y=y]=\sum_x x\Prob(X=x\mid Y=y),
\]
provided the sum is absolutely convergent.
\end{definition}


\begin{theorem}[Linearity of Expectation]
For integrable random variables $X, Y$ and scalars $a, b$,
\[
 \Exp[aX+bY]=a\Exp[X]+b\Exp[Y].
\]
No independence assumption is needed. If $X=\sum_{i=1}^n\indicator{A_i}$ counts how many of the events $A_i$ occur, then
\[
 \Exp[X]=\sum_{i=1}^n\Prob(A_i).
\]
\end{theorem}

\begin{formula}[Variance and Covariance]
For square-integrable real random variables $X, Y$,
\[
 \Var(X)=\Exp[X^2]-\Exp[X]^2,
\]
\[
 \Cov(X, Y)=\Exp[XY]-\Exp[X]\Exp[Y],
\]
\[
 \Var(X+Y)=\Var(X)+\Var(Y)+2\Cov(X, Y).
\]
If $X$ and $Y$ are independent, then $\Cov(X, Y)=0$.
\end{formula}

\begin{theorem}[Law of Total Expectation and Total Variance]
If $X$ is integrable and $Y$ is a random variable, then
\[
 \Exp[X]=\Exp[\Exp[X\mid Y]].
\]
If $X$ is square-integrable, then
\[
 \Var(X)=\Exp[\Var(X\mid Y)]
 +\Var(\Exp[X\mid Y]).
\]
\end{theorem}

\begin{theorem}[Laws of Large Numbers]
Let $X_1,X_2,\dots$ be independent identically distributed real random variables with $\Exp|X_1|<\infty$ and common mean $\mu$. Then
\[
 \frac1n\sum_{k=1}^n X_k\longrightarrow\mu\quad\text{a.s.},
\]
and hence also in probability (strong and weak laws of large numbers).
\end{theorem}

\begin{theorem}[Classical Central Limit Theorem]
Let $X_1,X_2,\dots$ be i.i.d. real random variables of common mean $\mu$ and variance $0<\sigma^2<\infty$. Then
\[
 \frac{\sum_{k=1}^n X_k-n\mu}{\sigma\sqrt n}
 \xrightarrow{\mathrm{d}}\Normal(0,1).
\]
\end{theorem}

\begin{theorem}[Tail-Sum Formula]
If $X$ is a nonnegative integer-valued random variable, then, allowing the value $+\infty$,
\[
 \Exp[X]=\sum_{k=1}^{\infty}\Prob(X\ge k).
\]
\end{theorem}

\begin{theorem}[Markov's Inequality]
If $X\ge0$ almost surely and $a>0$, then
\[
 \Prob(X\ge a)\le\frac{\Exp[X]}a,
\]
provided the right-hand side is finite.
\end{theorem}

\begin{theorem}[Chebyshev's Inequality]
If $X$ has finite variance and $t>0$, then
\[
 \Prob\!\left(|X-\Exp[X]|\ge t\right)
 \le\frac{\Var(X)}{t^2}.
\]
\end{theorem}

\begin{theorem}[Jensen's Inequality for Expectations]
For an integrable random variable $X$ taking values in an interval $I$ and a convex function $f\colon I\to\R$, whenever the relevant expectations exist,
\[
 f(\Exp[X])\le\Exp[f(X)].
\]
\end{theorem}

\begin{theorem}[Chernoff Bounds for Binomial Random Variables]
If $X\sim\Bin(n, p)$ has mean $\mu=np>0$, then for $\delta>0$,
\[
 \Prob(X\ge(1+\delta)\mu)
 \le\left(\frac{e^\delta}{(1+\delta)^{1+\delta}}\right)^\mu,
\]
and for $0<\delta<1$,
\[
 \Prob(X\le(1-\delta)\mu)\le e^{-\mu\delta^2/2}.
\]
\end{theorem}

\begin{theorem}[The First-Moment Method]
If $X$ is a nonnegative integer-valued random variable and $\Exp[X]<1$, then $\Prob(X=0)>0$. More generally, if $\Exp[X]>0$ and $X\ge0$, then $\Prob(X>0)>0$.
\end{theorem}

\begin{theorem}[Second-Moment Inequality]
For a real random variable $X$ with finite second moment and $\Exp[X^2]>0$,
\[
 \Prob(X\ne0)\ge\frac{(\Exp[X])^2}{\Exp[X^2]}.
\]
If $X\ge0$ almost surely, then for $0<\theta<1$ the Paley--Zygmund inequality gives
\[
 \Prob(X>\theta\Exp[X])
 \ge(1-\theta)^2\frac{(\Exp[X])^2}{\Exp[X^2]}.
\]
\end{theorem}

\begin{definition}[Filtration and Martingale]
A filtration $(\mathcal{F}_n)_{n\ge0}$ is an increasing sequence of sub-$\sigma$-algebras. An adapted process $(M_n)$ is a martingale if each $M_n$ is integrable and
\[
 \Exp[M_{n+1}\mid\mathcal{F}_n]=M_n\qquad\text{a.s.}
\]
A stopping time $\tau$ satisfies $\{\tau\le n\}\in\mathcal{F}_n$ for all $n$.
\end{definition}

\begin{theorem}[Optional Stopping Theorem, Bounded Case]
If $(M_n,\mathcal{F}_n)$ is a discrete-time martingale and $\tau$ is a bounded stopping time, then $M_\tau$ is integrable and
\[
 \Exp[M_\tau]=\Exp[M_0].
\]
\end{theorem}

\begin{formula}[Uniform Random Permutations]
If $\pi$ is uniformly distributed over the permutations of $[n]$, then the expected number of fixed points is $1$ for $n\ge1$, and the expected number of cycles is
\[
 H_n=\sum_{k=1}^n\frac1k.
\]
The probability that $\pi$ is a derangement is $D_n/n!$.
\end{formula}

\begin{formula}[Coupon Collector]
If one independently samples from $n$ equally likely types until all types have appeared, the expected number of samples is
\[
 nH_n=n\sum_{k=1}^n\frac1k.
\]
\end{formula}

\begin{formula}[Gambler's Ruin Probability]
For a nearest-neighbor random walk on $\{0, 1, \dots, N\}$ absorbed at $0, N$, with independent steps $+1$ of probability $p$ and $-1$ of probability $q=1-p$, the probability of reaching $N$ before $0$ from $i\in\{0, \dots, N\}$ is
\[
 \Prob_i(\text{hit }N\text{ before }0)=
 \begin{cases}
 \displaystyle\frac{1-(q/p)^i}{1-(q/p)^N}, &p\ne q, \;0<p<1, \\[0.7em]
 \displaystyle\frac{i}{N}, &p=q=\frac12.
 \end{cases}
\]
At $p=0$ or $p=1$, the probability is determined by the deterministic walk.
\end{formula}

\begin{theorem}[Lov\'asz Local Lemma, Symmetric Form]
Let $A_1, \dots, A_n$ be events with a dependency graph of maximum degree at most $d$: each $A_i$ is independent of the collection of all events corresponding to its nonneighbors. If $\Prob(A_i)\le p$ for all $i$ and
\[
 ep(d+1)\le1,
\]
then $\Prob\!\left(\bigcap_{i=1}^n A_i^c\right)>0$.
\end{theorem}
